% MinecrafTeX showcase: a short, PhD-level proof of the Central Limit Theorem,
% rendered entirely in Minecraft pixels.  Compile with:  lualatex clt.tex
\documentclass[11pt]{article}
\usepackage[a4paper,margin=2.3cm]{geometry}
\usepackage{amsmath}
\usepackage[scale=1.5]{minecraftex}

\setlength{\parindent}{0pt}
\setlength{\parskip}{6pt}

% Convergence in distribution: "d" over a right arrow (the arrow comes from the
% pixel font's U+2192 glyph, which \longrightarrow aliases to).
\newcommand{\dto}{\mathrel{\overset{\mathrm{d}}{\longrightarrow}}}
\newcommand{\E}{\mathbb{E}}
\newcommand{\R}{\mathbb{R}}
\newcommand{\N}{\mathbb{N}}
% End-of-proof box drawn with a TeX rule so it needs no extra font glyph.
\newcommand{\proofend}{\rule[0.05ex]{1.1ex}{1.1ex}}

\newcommand{\heading}[1]{\bigskip{\Large #1}\par\nobreak\smallskip}

\begin{document}

\begin{center}
{\huge The Central Limit Theorem}\\[4pt]
{\large A characteristic-function proof, in Minecraft pixels}
\end{center}

\heading{1.\quad Statement}

Let $X_1, X_2, \dots$ be independent, identically distributed random variables
on a probability space, with finite mean and variance
\[
  \mu = \E[X_1], \qquad \sigma^2 = \E\!\left[(X_1 - \mu)^2\right] \in (0, \infty).
\]
Write the partial sum $S_n = \sum_{k=1}^{n} X_k$ and its standardisation
\[
  Z_n = \frac{S_n - n\mu}{\sigma \sqrt{n}}.
\]

\textbf{Theorem (CLT).}\; As $n \to \infty$,
\[
  Z_n \;\dto\; Z, \qquad Z \sim \mathcal{N}(0, 1),
\]
that is, for every $z \in \R$,
\[
  \Pr(Z_n \leq z) \;\longrightarrow\;
  \Phi(z) = \int_{-\infty}^{z} \frac{1}{\sqrt{2\pi}}\, e^{-x^2/2}\, dx.
\]

\heading{2.\quad The characteristic function}

For a random variable $X$ the characteristic function is
\[
  \varphi_X(t) = \E\!\left[ e^{itX} \right]
              = \int_{-\infty}^{\infty} e^{itx}\, dF_X(x), \qquad t \in \R.
\]
Two facts drive the proof. First, independence factorises it: if $X$ and $Y$ are
independent then $\varphi_{X+Y} = \varphi_X \cdot \varphi_Y$. Second, two moments
give a quadratic expansion at the origin,
\[
  \varphi_X(t) = 1 + it\,\E[X] - \frac{t^2}{2}\,\E[X^2] + o(t^2)
  \qquad (t \to 0).
\]

\newpage

\heading{3.\quad Proof}

By replacing $X_k$ with $(X_k - \mu)/\sigma$ we may assume $\mu = 0$ and
$\sigma^2 = 1$, so that $\E[X_1] = 0$, $\E[X_1^2] = 1$, and
\[
  Z_n = \frac{1}{\sqrt{n}} \sum_{k=1}^{n} X_k.
\]
Fix $t \in \R$. Using independence and the identical distribution of the $X_k$,
\[
  \varphi_{Z_n}(t)
    = \E\!\left[ \exp\!\left( \frac{it}{\sqrt{n}} \sum_{k=1}^{n} X_k \right) \right]
    = \prod_{k=1}^{n} \varphi_{X_1}\!\left( \frac{t}{\sqrt{n}} \right)
    = \left[ \varphi_{X_1}\!\left( \frac{t}{\sqrt{n}} \right) \right]^{n}.
\]
Since $\E[X_1] = 0$ and $\E[X_1^2] = 1$, the expansion of Section~2 at the point
$t/\sqrt{n} \to 0$ gives
\[
  \varphi_{X_1}\!\left( \frac{t}{\sqrt{n}} \right)
    = 1 - \frac{t^2}{2n} + o\!\left( \frac{1}{n} \right).
\]
Therefore, taking logarithms and using $\log(1 + u) = u + o(u)$,
\[
  \log \varphi_{Z_n}(t)
    = n \log\!\left( 1 - \frac{t^2}{2n} + o\!\left(\tfrac{1}{n}\right) \right)
    = n \left( -\frac{t^2}{2n} + o\!\left(\tfrac{1}{n}\right) \right)
    \;\longrightarrow\; -\frac{t^2}{2}.
\]
Equivalently, with the classical limit $\bigl(1 + \tfrac{a}{n}\bigr)^{n} \to e^{a}$,
\[
  \varphi_{Z_n}(t) \;\longrightarrow\; e^{-t^2/2}
  \qquad \text{for every } t \in \R.
\]
The right-hand side is exactly the characteristic function of the standard normal
law, since
\[
  \int_{-\infty}^{\infty} e^{itx}\, \frac{1}{\sqrt{2\pi}}\, e^{-x^2/2}\, dx
    = e^{-t^2/2}.
\]
By Levy's continuity theorem, pointwise convergence of characteristic functions
to a function continuous at $0$ implies convergence in distribution. Hence
\[
  Z_n \;\dto\; Z \sim \mathcal{N}(0, 1). \qquad \proofend
\]

\heading{4.\quad Remark}

The same argument, expanded one order further, exhibits the speed of convergence:
if $\E\,|X_1|^3 < \infty$ then for each fixed $t$,
\[
  \left| \varphi_{Z_n}(t) - e^{-t^2/2} \right|
    \;\leq\; \frac{C\,|t|^3}{\sqrt{n}},
\]
which is the analytic heart of the Berry--Esseen bound on
$\sup_{z \in \R} \bigl| \Pr(Z_n \leq z) - \Phi(z) \bigr|$.

\end{document}
